% !PRINTED 297
\weberpage{297}
\begin{center}
\S.\,75.\\
\gesperrt{Metacyklische Gleichungen, deren Grad eine Primzahlpotenz ist.}
\end{center}

Nach den letzten S\"atzen concentrirt sich das Interesse weiterer Untersuchungen \"uber metacyklische Gleichungen haupts\"achlich auf den Fall, dass der Grad der Gleichung eine Potenz $p^k$ einer Primzahl $p$ ist. Wir setzen die Gleichung als irreducibel voraus und beschr\"anken uns auf die Betrachtung \gesperrt{primitiver Gleichungen}; denn die Aufl\"osung der imprimitiven reducirt sich auf die successive L\"osung zweier (oder mehrerer) primitiver Gleichungen, deren Grade gleichfalls Potenzen von $p$ sind, und die, wenn die urspr\"ungliche Gleichung metacyklisch ist, auch metacyklisch sein m\"ussen.

Wir nehmen also jetzt an, es sei $P$ die Gruppe einer irreduciblen primitiven metacyklischen Gleichung $f(x) = 0$ vom Grade $p^k$. Nach Bd. I, \S.\,158, 2. muss nicht nur $P$ selbst, sondern alle seine Normaltheiler (mit Ausnahme der Einheitsgruppe) noch \gesperrt{transitiv} sein. Nun wenden wir den in \S.\,8 allgemein bewiesenen Satz IV. an, nach dem $P$ einen Normaltheiler $Q$ mit lauter commutativen Elementen besitzt, und wenn mehrere solche Theiler vorhanden sind, so verstehen wir unter $Q$ einen von ihnen, dessen Grad m\"oglichst niedrig, aber noch gr\"osser als 1 ist. Diese Gruppe $Q$ ist also gleichfalls noch transitiv. $Q$ kann aber keinen von der Einheit verschiedenen echten Theiler mehr haben, der zugleich Normaltheiler von $P$ ist, weil sonst dieser an die Stelle von $Q$ treten w\"urde.

Wenn durch die geh\"origen Adjunctionen die Gruppe unserer Gleichung von $P$ auf $Q$ reducirt ist, so ist $f(x) = 0$ in dem erweiterten Rationalit\"atsbereiche zu einer \gesperrt{Abel'schen Gleichung} geworden, und da sie noch irreducibel geblieben ist, so ist nach Bd. I, \S.\,162 der Grad der Gleichung gleich dem Grade der Gruppe; d.\,h. $Q$ ist eine \gesperrt{Abel'sche Gruppe} vom Grade $p^k$. Die Grade aller Elemente von $Q$ sind also Potenzen von $p$. Wir wollen nachweisen, dass ausser dem Einheitselemente in $Q$ \gesperrt{nur Elemente vom Grade $p$ selbst vorkommen.}

Nehmen wir an, es sei $p^\lambda$ der h\"ochste Grad, der unter den Elementen von $Q$ vorkommt, und es sei $\lambda > 1$. Dann bilden alle Elemente von $Q$, deren Grad ein Theiler von $p^{\lambda-1}$ ist, einen
%CONT>>
